\chapter{Introdução}
\label{cap:introducao}

\section{Contextualização e Motivação}
\label{sec:contextualizacao}

Os sistemas de controle em malha fechada constituem a base operacional de grande parte da infraestrutura tecnológica e industrial moderna. Desde processos químicos contínuos e plantas de geração de energia até sistemas automotivos avançados e veículos aéreos não tripulados, a capacidade de manter variáveis de processo reguladas em torno de referências pré-determinadas, mesmo sob a atuação de perturbações externas e incertezas paramétricas, é indispensável para garantir eficiência, segurança e viabilidade operacional \cite{ogata2010, dorf2017}.

Nesse cenário, o algoritmo Proporcional-Integral-Derivativo (PID) consolida-se historicamente como a estratégia de controle por realimentação mais empregada em aplicações industriais, respondendo por mais de 90\% das malhas de controle ativas no setor produtivo \cite{astrom2006, nise2020}. A ampla aceitação do controlador PID decorre de sua formulação conceitualmente intuitiva --- baseada no tratamento do erro passado (integral), presente (proporcional) e futuro estimado (derivativo) --- aliada à sua robustez estrutural e à independência de modelos matemáticos complexos de ordem elevada para sua implementação preliminar.

Não obstante sua ubiquidade e relevância prática, o ensino e a aprendizagem de sistemas de controle na formação de engenheiros --- particularmente em cursos de Engenharia de Computação --- enfrentam desafios pedagógicos expressivos. Tradicionalmente, as ementas curriculares enfatizam o domínio analítico no domínio da frequência e no plano complexo (Transformadas de Laplace e Fourier, Diagramas de Bode e Nyquist, Lugar Geométrico das Raízes). Embora esse embasamento formal seja estritamente necessário, observa-se com frequência um distanciamento substancial entre a representação matemática abstrata e a dinâmica física concreta do sistema real \cite{dormido2002}.

Muitas vezes, o estudante manipula funções de transferência sem compreender visualmente como o ajuste do ganho derivativo $K_d$ amortece o movimento de um veículo, de que forma o ganho integral $K_i$ atua para vencer o atrito estático em um aclive, ou quais são as consequências severas da saturação do atuador associada ao fenômeno de \textit{integrator windup}. Essa lacuna epistemológica entre a teoria e a intuição física contribui para índices expressivos de desmotivação e evasão em disciplinas fundamentais da área de automação e controle.

\section{Problema de Pesquisa}
\label{sec:problema_pesquisa}

A consolidação prática dos conceitos de engenharia de controle tradicionalmente depende de duas vias: a experimentação em bancadas didáticas físicas ou o uso de simuladores computacionais. Entretanto, ambas as abordagens apresentam restrições relevantes no contexto universitário contemporâneo:

\begin{enumerate}[label=\alph*)]
    \item \textbf{Custo e Manutenção de Bancadas Físicas:} Laboratórios físicos com servomecanismos, plantas de nível ou túneis de vento exigem investimentos financeiros expressivos em aquisição, calibração periódica e manutenção de atuadores e sensores, inviabilizando a disponibilização de estações individuais para turmas numerosas;
    \item \textbf{Barreiras de Licenciamento de Softwares Proprietários:} Ferramentas industriais consagradas, a exemplo do MATLAB/Simulink e do LabVIEW, impõem custos proibitivos de licenciamento para discentes, exigem computadores com capacidade expressiva de processamento e demandam uma curva de aprendizado substancial apenas para a operação do ambiente de software;
    \item \textbf{Limitações de Simuladores Acadêmicos Tradicionais:} A maioria das alternativas computacionais gratuitas disponíveis na web limita-se a gráficos bidimensionais estáticos e controles deslizantes parametrizados, sem permitir que o estudante programe diretamente o código do algoritmo em malha fechada, nem que observe a resposta física da planta em uma cena tridimensional realista com perturbações dinâmicas de relevo e atrito.
\end{enumerate}

Diante deste panorama, formula-se o problema de pesquisa que orienta o presente trabalho:

\begin{quote}
\textit{Como conceber e implementar uma arquitetura computacional baseada na web que combine simulação numérica em tempo real, computação gráfica tridimensional e decomposição analítica do esforço de controle para mitigar a abstração conceitual e os altos custos de infraestrutura no ensino e na parametrização de controladores PID?}
\end{quote}

\section{Hipótese de Pesquisa}
\label{sec:hipotese}

A hipótese científica deste trabalho estabelece que a disponibilização de uma plataforma web interativa, dotada de renderização física tridimensional em tempo real sincronizada a um osciloscópio digital com decomposição explícita das forças proporcional ($u_P$), integral ($u_I$) e derivativa ($u_D$), acoplada a um módulo assistido de sintonia em malha aberta (Ziegler-Nichols), proporciona uma assimilação conceitual superior e mais intuitiva da dinâmica transitória e de regime permanente do que métodos restritos ao equacionamento puramente analítico ou a simuladores bidimensionais estáticos.

\section{Objetivos}
\label{sec:objetivos}

\subsection{Objetivo Geral}
\label{subsec:obj_geral}

Projetar, implementar e validar uma plataforma web interativa para simulação e sintonia de controladores PID em malha fechada, integrando renderização física tridimensional em tempo real, análise de telemetria didática e geração automatizada de relatórios técnicos de desempenho para apoio ao ensino de engenharia.

\subsection{Objetivos Específicos}
\label{subsec:obj_especificos}

Para viabilizar a consecução do objetivo geral, estabelecem-se as seguintes metas específicas de desenvolvimento e pesquisa:

\begin{enumerate}
    \item \textbf{Modelar matematicamente} a dinâmica física e cinemática de plantas representativas (veículo terrestre sob atrito de rolamento, arrasto aerodinâmico e componentes gravitacionais em relevos inclinados com $\alpha(x)$ variável);
    \item \textbf{Desenvolver a arquitetura computacional frontend} em WebGL/Three.js e JavaScript vanilla, incorporando um editor de código isolado em \textit{sandbox} capaz de executar algoritmos do usuário com determinismo temporal a taxas de atualização fluidas ($\ge 30\text{ FPS}$);
    \item \textbf{Implementar a instrumentação e osciloscópio digital} multicanal com capacidade de monitoramento síncrono da variável de processo $y(t)$, setpoint $r(t)$, decomposição individual das parcelas de controle ($u_P, u_I, u_D$) e sinalização de saturação do atuador;
    \item \textbf{Construir um assistente automatizado de sintonia clássica} em malha aberta, aplicando o método da reta tangente à curva de reação para identificação de modelos FOPDT ($K, \tau, \theta$) e cálculo de parâmetros de sintonia baseados nos critérios de Ziegler-Nichols;
    \item \textbf{Projetar a infraestrutura de backend} em Python (FastAPI) conteinerizada via Docker, provendo serviços de persistência de ensaios, exportação de conjuntos de dados em formato CSV e emissão automática de relatórios técnicos em PDF com diagnósticos de controle;
    \item \textbf{Validar experimentalmente} a precisão numérica da simulação frente a soluções analíticas exatas calculadas por métodos numéricos consolidados, quantificando métricas de resposta ao degrau ($t_r, t_s, M_p, e_{ss}$) e o erro médio quadrático ($\text{RMSE}$).
\end{enumerate}

\section{Justificativa e Relevância}
\label{sec:justificativa}

A realização deste trabalho justifica-se sob quatro perspectivas complementares:

\begin{itemize}
    \item \textbf{Relevância Pedagógica:} O ambiente oferece feedback visual imediato, permitindo que o estudante correlacione causa e efeito em tempo real. A possibilidade de visualizar as forças do atuador de forma decomposta e de presenciar visualmente o veículo respondendo a aclives e declives promove a fixação dos conceitos fundamentais de estabilidade e transitório;
    \item \textbf{Acessibilidade e Inclusão Tecnológica:} Por ser executada integralmente no navegador web (\textit{browser-based}) e distribuída sob arquitetura aberta, a ferramenta independe de estações de trabalho de alto custo ou da aquisição de licenças proprietárias, democratizando o acesso a laboratórios práticos na Universidade Federal de Alagoas (UFAL) e em outras instituições de ensino;
    \item \textbf{Desafio Técnico em Engenharia de Computação:} O desenvolvimento congrega áreas essenciais do núcleo profissional do curso, abrangendo engenharia de software distribuída (FastAPI, Docker), computação gráfica e matemática aplicada (Three.js, WebGL), integração numérica determinística e segurança computacional em ambientes de execução de código não confiável (\textit{client-side sandbox});
    \item \textbf{Apoio à Pesquisa e Experimentação Científica:} A exportação de dados estruturados em CSV e a emissão de relatórios automatizados conferem rigor científico aos ensaios, permitindo que a plataforma seja utilizada não apenas no ensino de graduação, mas também como bancada preliminar de testes para prototipagem de algoritmos de controle.
\end{itemize}

\section{Organização do Trabalho}
\label{sec:organizacao}

A presente monografia encontra-se estruturada em cinco capítulos, distribuídos conforme o roteiro a seguir:

\begin{itemize}
    \item O \textbf{Capítulo \ref{cap:introducao} (Introdução)} contextualiza o tema, delimita o problema de pesquisa, explicita as hipóteses e estabelece os objetivos e a relevância do estudo;
    \item O \textbf{Capítulo 2 (Fundamentação Teórica e Trabalhos Correlatos)} revisita os conceitos de controle clássico por realimentação PID, métricas de desempenho transitório, fenômenos de saturação, métodos de identificação de sistemas (FOPDT e Ziegler-Nichols), além de avaliar o estado da arte dos laboratórios virtuais de ensino;
    \item O \textbf{Capítulo 3 (Desenvolvimento e Arquitetura do Sistema)} expõe a metodologia do projeto, a modelagem física e matemática das plantas, a arquitetura de software (frontend e backend), a concepção do osciloscópio multicanal, o módulo de sintonia assistida e o ambiente de conteinerização;
    \item O \textbf{Capítulo 4 (Resultados Experimentais e Discussão)} apresenta os ensaios de validação numérica frente a curvas analíticas teóricas, a análise de desempenho do algoritmo PID sob perturbações de relevo, a eficácia do módulo de sintonia e o perfil de desempenho computacional da aplicação;
    \item O \textbf{Capítulo 5 (Conclusão e Trabalhos Futuros)} sintetiza as principais contribuições alcançadas, relata as limitações identificadas e propõe direcionamentos para trabalhos subsequentes.
\end{itemize}
